\chapter{Злоякісний програмний код як об'єкт протидії: поняття, таксономія та простір загроз}
\chaptershort{Злоякісний програмний код як об'єкт протидії}
\label{chap:ch01}

Злоякісний програмний код став одним із головних інструментів у сучасних кібератаках і водночас одним із найскладніших об'єктів захисту. На відміну від пасивних загроз, він є активним, самодостатнім артефактом, що виконується на обчислювальній системі й цілеспрямовано порушує її безпеку: викрадає дані, шифрує або знищує їх, перехоплює керування, забезпечує прихований доступ зловмисника. За останні десятиліття злоякісний код пройшов шлях від експериментальних саморозмножуваних програм до промислової екосистеми з розподілом праці, ринком послуг і державним спонсоруванням, а кількість нових зразків зросла до сотень мільйонів на рік. Для України ця проблематика має особливу гостроту: країна стала полігоном найруйнівніших кампаній в історії (NotPetya, Industroyer, хвиля деструктивних «вайперів» 2022~року), що зробило протидію злоякісному коду питанням національної безпеки. Розділ формує понятійний апарат і системне бачення предмета: дає означення й термінологію, будує таксономію злоякісного коду, простежує його еволюцію, описує життєвий цикл злоякісної кампанії, модель порушника та простір загроз, аналізує український контекст і окреслює принципи протидії й структуру курсу.

\section{Поняття та термінологія}\label{sec:ch01-terms}

\subsection{Означення злоякісного програмного коду}

\term{Злоякісний програмний код} (ЗПК; у міжнародній літературі --- \emph{malware}, від \emph{malicious software}) --- це програмний код, навмисно створений для виконання несанкціонованих, шкідливих для власника інформаційної системи дій. Ключовими в означенні є три ознаки: \emph{навмисність} (код створено з наміром завдати шкоди, що відрізняє ЗПК від звичайних програмних помилок), \emph{несанкціонованість} (дії виконуються без відома та згоди законного власника) і \emph{функціональна спрямованість на порушення безпеки} (конфіденційності, цілісності, доступності, підзвітності ресурсів). У вітчизняній нормативній традиції (НД~ТЗІ~1.1-003-99~\cite{ndtzi-1-1-003}) споріднене поняття --- «програмна закладка» та «комп'ютерний вірус»; у Кримінальному кодексі України~\cite{kk-ukr} вживається термін «шкідливі програмні чи технічні засоби». У цьому посібнику терміни «злоякісний програмний код» і «шкідливе програмне забезпечення» (ШПЗ) вживаються як синоніми.

Принципово, що «злоякісність» визначається не внутрішніми властивостями коду, а \emph{наміром} і \emph{контекстом} його використання. Той самий технічний прийом (наприклад, віддалене виконання команд, шифрування файлів, перехоплення введення) може бути легітимною функцією адміністрування чи захисту або складовою атаки. Ця двоїстість (\term{dual-use}) є фундаментальною для предметної галузі: вона ускладнює автоматичне розрізнення й зумовлює принципову неможливість суто технічного визначення злоякісності у загальному випадку (розділ~\ref{chap:ch03}).

Прикладом служить клас інструментів подвійного призначення: засоби віддаленого адміністрування (RAT у легітимному сенсі), утиліти стиснення й шифрування, фреймворки тестування на проникнення (Cobalt~Strike, Metasploit), системні утиліти командного рядка. Усі вони мають законне застосування, але масово зловживаються в атаках; розмежування спирається на поведінковий контекст (хто, коли, за яких обставин запустив, з якими параметрами), а не на сам факт наявності коду. Саме тому сучасні засоби захисту оцінюють не лише артефакт, а й ланцюг подій навколо його виконання (розділи~\ref{chap:ch07}, \ref{chap:ch10}). Окремою складністю є \term{потенційно небажані програми} (PUP/PUA) --- рекламне ПЗ, агресивні тулбари, програми стеження, що формально не є злоякісними, але порушують інтереси користувача; їх класифікація залежить від політики й згоди користувача.

\subsection{Злоякісний код у системі понять безпеки}

Злоякісний код є \emph{засобом реалізації загрози}, а не самою загрозою. У термінах управління ризиками (розділ~\ref{chap:ch04}) джерело загрози (зловмисник) використовує ЗПК для експлуатації вразливості з метою завдати впливу на актив. Протидія злоякісному коду (\term{anti-malware}, захист від шкідливого коду) --- це сукупність технологій, процесів і організаційних заходів, спрямованих на запобігання проникненню, виявлення, аналіз, стримування, видалення ЗПК і відновлення після його дії. Відповідно до функцій рамки NIST~CSF~2.0~\cite{nist-csf2} протидія охоплює ідентифікацію, захист, виявлення, реагування й відновлення, а настанова NIST~SP~800-83~\cite{nist-83} присвячена саме запобіганню інцидентам зі шкідливим кодом та їх обробці.

\section{Таксономія злоякісного програмного коду}\label{sec:ch01-taxonomy}

Злоякісний код класифікують за кількома незалежними ознаками~\cite{ch01:grgof,aycock,szor}: за \emph{механізмом поширення й автономності}, за \emph{функціональним призначенням} (корисним навантаженням) і за \emph{способом маскування}. Ці осі ортогональні: конкретний зразок одночасно належить до класів за кожною з них (наприклад, самопоширюваний черв'як-вимагач з руткіт-компонентом).

\subsection{Класифікація за механізмом поширення}

За здатністю до самопоширення й потребою в носії розрізняють три базові типи. \term{Вірус} --- код, що вбудовується в інший файл-носій (виконуваний файл, документ із макросом, завантажувальний сектор) і розмножується, заражаючи нові носії під час виконання. \term{Черв'як} (worm) --- самостійна програма, що поширюється мережею без носія й без участі користувача, експлуатуючи вразливості або слабкі облікові дані. \term{Троян} (троянський кінь) --- програма, що маскується під легітимну й не має механізму самопоширення, покладаючись на дії користувача (запуск, встановлення). Формальну основу теорії вірусів заклав Коен~\cite{ch01:cohen,ch01:cohen-thesis}, спираючись на модель саморозмножуваних автоматів фон Неймана~\cite{ch01:vonneumann}; перший масштабний мережевий черв'як проаналізував Споффорд~\cite{ch01:spafford}.

Коен також довів принциповий результат: задача точного виявлення вірусів у загальному випадку є алгоритмічно нерозв'язною (зводиться до проблеми зупинки), що має прямі наслідки для теорії протидії (розділ~\ref{chap:ch03}). Механізм самопоширення визначає \emph{динаміку епідемії}: для черв'яків, що сканують адресний простір, початкове зростання числа заражених вузлів є експоненційним і описується епідемічними моделями (розділ~\ref{chap:ch04}), чим пояснюється швидкість історичних епідемій на кшталт Slammer, який уразив більшість вразливих вузлів за лічені хвилини. Швидкість поширення обмежується щільністю вразливих вузлів і пропускною здатністю, а не можливостями самого коду, тому своєчасне усунення вразливостей (розділ~\ref{chap:ch10}) кардинально змінює динаміку.

\subsection{Класифікація за функціональним призначенням}

За корисним навантаженням виділяють (табл.~\thechapter.1): \term{бекдор} (backdoor) --- забезпечує прихований віддалений доступ; \term{завантажувач} (downloader/dropper) --- доставляє й встановлює інші компоненти; \term{програма-вимагач} (ransomware) --- шифрує дані з вимогою викупу (розділ~\ref{chap:ch12}); \term{вайпер} (wiper) --- незворотно знищує дані, імітуючи вимагач або без маскування; \term{шпигунське ПЗ} (spyware) --- викрадає інформацію; \term{клавіатурний шпигун} (keylogger) --- перехоплює введення; \term{банківський троян} --- краде фінансові облікові дані; \term{руткіт} (rootkit) --- приховує присутність інших компонентів, модифікуючи систему; \term{бот} --- робить вузол частиною керованої мережі (ботнету); \term{криптомайнер} --- використовує ресурси для видобутку криптовалюти.

\tabhead{Таблиця~\thechapter.1 --- Класифікація злоякісного коду за функціональним призначенням}

\begin{center}\footnotesize
\begin{tabularx}{\linewidth}{L{24mm} Y L{16mm}}
\toprule
Клас & Призначення (корисне навантаження) & Приклад \\
\midrule
Бекдор & прихований віддалений доступ і керування & багато APT-імплантів \\
Завантажувач & доставка й встановлення інших компонентів & Emotet, dropper'и \\
Вимагач & шифрування даних із вимогою викупу & розділ~\ref{chap:ch12} \\
Вайпер & незворотне знищення даних & HermeticWiper, CaddyWiper \\
Шпигунське ПЗ & викрадення інформації, стеження & інфостілери \\
Банківський троян & крадіжка фінансових даних & розділ~\ref{chap:ch13} \\
Руткіт & приховування присутності в системі & розділ~\ref{chap:ch05} \\
Бот & включення вузла до ботнету & розділ~\ref{chap:ch11} \\
Промислове ЗПК & вплив на технологічні процеси & Industroyer, Stuxnet \\
\bottomrule
\end{tabularx}
\end{center}

\subsection{Класифікація за способом маскування}

За здатністю ухилятися від аналізу й виявлення розрізняють (розділ~\ref{chap:ch08}): \term{пакований} код (стиснутий/зашифрований пакувальником, що розгортається в пам'яті); \term{поліморфний} (змінює власний код у кожному поколінні, зберігаючи функцію, через шифрування зі змінним ключем); \term{метаморфний} (переписує власне тіло без шифрування); \term{безфайловий} (виконується лише в пам'яті, не залишаючи файлу на диску, розділ~\ref{chap:ch13}); код із \term{антианалізом} (виявляє пісочницю, налагоджувач, віртуальне середовище й змінює поведінку).

\subsection{Багатовимірність класифікації}

Жоден зразок не вичерпується одним класом: сучасний ЗПК є \emph{модульним} і \emph{багатофункціональним}. Тому практичні таксономії є фасетними (багатовимірними), що відображають одночасну належність зразка до класів за різними осями; розроблення такої багатовимірної фасетної таксономії, зокрема технік обфускації, є предметом окремих досліджень~\cite{ch01:shevchenko}. Для узгодженого опису \emph{поведінки} (а не формальної належності) галузевим стандартом де-факто є база знань MITRE~ATT\&CK~\cite{mitre-attack}, що каталогізує тактики й техніки, а база MITRE~D3FEND~\cite{mitre-defend} --- відповідні контрзаходи (розділ~\ref{chap:ch04}).

\section{Еволюція злоякісного коду}\label{sec:ch01-evolution}

Еволюцію ЗПК зручно подати як зміну \emph{домінівної мотивації} та \emph{моделі організації} (рис.~\thechapter.1). \emph{Епоха цікавості} (1980-ті -- середина 1990-х): експериментальні віруси й черв'яки, що поширювалися через дискети й ранні мережі; мотив --- дослідницький чи вандальний. \emph{Епоха масових епідемій} (кінець 1990-х -- 2000-ті): поштові й мережеві черв'яки (ILOVEYOU, Code~Red, Slammer), що уражали мільйони вузлів за години; поява комерційної мотивації (спам-ботнети, банківські трояни). \emph{Епоха кіберзлочинності як індустрії} (2010-ті): професіоналізація, ринок шкідливого коду як послуги (\term{MaaS}), партнерські програми, спеціалізація (розробники, оператори, брокери доступу). \emph{Епоха державних операцій та деструкції} (з початку 2010-х): кіберзброя проти промислових систем (Stuxnet~\cite{stuxnet,ch01:zetter}), складні постійні загрози (APT), деструктивні кампанії проти критичної інфраструктури й цілих держав (Україна, 2015--2022~рр.~\cite{ch01:greenberg}).

\begin{tikzfig}
\begin{tikzpicture}[font=\scriptsize,
  b/.style={draw,rounded corners=2pt,minimum height=7mm,text width=24mm,align=center},
  >={Stealth[length=1.6mm]}]
\node[b] (e1) at (0,0) {1980-ті\\цікавість,\\віруси};
\node[b] (e2) at (0,-1.1) {1990--2000-ті\\масові епідемії,\\черв'яки};
\node[b] (e3) at (0,-2.2) {2010-ті\\індустрія (MaaS),\\кіберзлочинність};
\node[b] (e4) at (0,-3.3) {2010-ті--тепер\\державні операції,\\деструкція, APT};
\draw[->] (e1) -- (e2); \draw[->] (e2) -- (e3); \draw[->] (e3) -- (e4);
\node[anchor=west,align=left] at (3,-1.1) {зростання: масштабу,};
\node[anchor=west,align=left] at (3,-1.5) {автоматизації, мотивації,};
\node[anchor=west,align=left] at (3,-1.9) {організованості й шкоди};
\draw[->] (2.6,-0.2) -- (2.6,-3.1);
\end{tikzpicture}
\figcap{Рис.~\thechapter.1 --- Еволюція домінівної мотивації та організації злоякісного коду}
\end{tikzfig}

Ключовою рисою сучасної епохи є \term{комодитизація} злоякісного коду --- перетворення його на товар і послугу. Склалася багаторівнева тіньова економіка з розподілом праці: розробники створюють шкідливий код і продають його або здають в оренду (\term{Malware-as-a-Service}); брокери початкового доступу (\term{Initial Access Brokers}) продають вже отриманий доступ до скомпрометованих мереж; оператори вимагачів працюють за партнерськими програмами (\term{Ransomware-as-a-Service}), ділячи викуп; окремо діють сервіси обфускації, тестування на виявлюваність антивірусами, оренди ботнетів і відмивання коштів. Ця спеціалізація різко знизила поріг входу: атаку здатен провести виконавець без глибоких технічних навичок, орендувавши готові інструменти. Для захисту це означає зростання кількості й різноманітності атак та розмивання межі між «простими» й «складними» загрозами, оскільки навіть малокваліфікований зловмисник отримує доступ до професійно розробленого коду.

Кількісно масштаб проблеми ілюструє зростання числа унікальних зразків. Якщо позначити через $N(t)$ накопичену кількість відомих зразків на момент $t$, то емпірично впродовж тривалого періоду вона зростала приблизно експоненційно:

\begin{equation}
N(t) = N_0\,e^{\,\alpha (t-t_0)},
\label{eq:ch01-growth}
\end{equation}

де $\alpha$ --- коефіцієнт зростання. За даними антивірусних лабораторій~\cite{ch01:avtest} щорічно реєструють сотні мільйонів нових зразків, а загальна кількість перевищила мільярд. Проста наслідкова оцінка пояснює кризу суто сигнатурного підходу: якщо щодня з'являється $\Delta N$ нових зразків, а лабораторія здатна створити $S$ сигнатур за добу, то частка невкритих зразків не менша за $1-S/\Delta N$; за $\Delta N \gg S$ сигнатурне виявлення принципово відстає (розділ~\ref{chap:ch06}), що й зумовило перехід до поведінкових і статистичних методів (розділи~\ref{chap:ch07}, \ref{chap:ch09}).

\section{Життєвий цикл злоякісної кампанії}\label{sec:ch01-lifecycle}

Злоякісний код є лише інструментом у ширшій \emph{операції}. Її життєвий цикл описують моделлю ланцюга ураження (\term{kill chain})~\cite{lockheed-kc}, детально розглянутою в розділі~\ref{chap:ch04}. Узагальнено кампанія проходить етапи: розвідка (збір інформації про ціль); озброєння (підготовка ЗПК і доставки); доставка (фішинг, заражені носії, вразливий сервіс, ланцюг постачання); експлуатація (виконання коду через вразливість чи дію користувача); встановлення (закріплення в системі, персистентність); командування й керування (зв'язок з інфраструктурою зловмисника, \term{C2}); дії за цілями (крадіжка, шифрування, знищення, рух мережею). Кожен етап є точкою потенційного виявлення й розриву ланцюга --- ключова ідея інтелектуально-керованого захисту: достатньо надійно перервати \emph{будь-який} обов'язковий етап, щоб зірвати операцію.

Життєвий цикл власне злоякісного артефакта вкладений у цей ширший цикл і охоплює: доставку й первинне виконання; розгортання (розпакування, розшифрування тіла в пам'яті --- розділ~\ref{chap:ch08}); перевірку середовища (антианаліз, виявлення пісочниці); закріплення (персистентність через автозапуск, служби, заплановані завдання, модифікацію системи --- розділ~\ref{chap:ch05}); підвищення привілеїв; встановлення зв'язку з C2; виконання корисного навантаження; за потреби --- самовидалення й приховування слідів. Розуміння цього циклу є основою як для аналізу (розділи~\ref{chap:ch06}, \ref{chap:ch07}), так і для виявлення: кожна фаза залишає характерні артефакти й поведінкові ознаки, що стають індикаторами компрометації. Важлива асиметрія часу: зловмиснику вигідно діяти швидко до спрацювання навантаження, тоді як захиснику критично скоротити час від проникнення до виявлення (\term{dwell time}), який для складних загроз історично вимірювався тижнями й місяцями (розділ~\ref{chap:ch13}).

\section{Модель порушника та мотивація}\label{sec:ch01-adversary}

Протидія має відповідати моделі порушника (розділ~\ref{chap:ch04}). Профілі порушників, що застосовують ЗПК, суттєво різняться за ресурсами, мотивацією й наполегливістю (табл.~\thechapter.2).

\tabhead{Таблиця~\thechapter.2 --- Профілі порушників, що застосовують злоякісний код}

\begin{center}\footnotesize
\begin{tabularx}{\linewidth}{L{26mm} L{22mm} Y}
\toprule
Профіль & Мотивація & Характер і ресурси \\
\midrule
Кіберзлочинці & фінансова вигода & масовість, MaaS, вимагачі, інфостілери; помірні ресурси \\
Організовані угруповання & великий прибуток & партнерські програми, подвійне вимагання; значні ресурси \\
Хактивісти & ідеологія, розголос & дефейси, DDoS, витоки; змінні ресурси \\
Інсайдери & помста, вигода & легітимний доступ; складно виявити \\
Державні APT & розвідка, вплив, деструкція & тривала прихована присутність, 0-day, кіберзброя; великі ресурси \\
\bottomrule
\end{tabularx}
\end{center}

Для кількісного порівняння корисна інтегральна оцінка \emph{рівня загрози}, що агрегує можливості й мотивацію порушника:

\begin{equation}
T = \frac{C \cdot M \cdot O}{D},
\label{eq:ch01-threat}
\end{equation}

де $C$ --- технічні можливості (capability); $M$ --- мотивація; $O$ --- наполегливість (persistence); $D$ --- стримувальні чинники (ризик викриття й покарання). Формула~\eqref{eq:ch01-threat} якісно пояснює особливу небезпеку державних APT: висока $C$ і $M$, екстремальна $O$ за практично відсутнього $D$ (атакувальник діє з юрисдикції, що його не переслідує). Для таких порушників традиційні засоби стримування не працюють, і захист має будуватися на припущенні компрометації (розділ~\ref{chap:ch13}).

\section{Простір загроз і поверхня атаки}\label{sec:ch01-surface}

\term{Поверхня атаки} для проникнення ЗПК --- сукупність шляхів доставки й виконання коду. Основні вектори: електронна пошта (фішинг з вкладеннями та посиланнями --- домінівний вектор за даними~\cite{ch01:dbir}); вебзагрози (заражені сайти, шкідлива реклама, завантаження); зовнішні носії (USB); експлуатація вразливих сервісів, доступних з мережі; компрометація облікових даних і легітимного віддаленого доступу; ланцюг постачання програмного забезпечення (розділ~\ref{chap:ch15}); мобільні застосунки з ненадійних джерел (розділ~\ref{chap:ch14}). Поверхня атаки зростає з ускладненням середовищ: хмарні, контейнерні, IoT- та OT-системи додають власні вектори (розділ~\ref{chap:ch14}).

Зменшення поверхні атаки --- перший і найдешевший рубіж протидії: обмеження вкладень і макросів, контроль застосунків, вимкнення непотрібних сервісів, сегментація, своєчасне оновлення (розділ~\ref{chap:ch10}). Проте жодне скорочення поверхні не усуває ризик повністю, тому необхідні глибоко ешелоновані виявлення й реагування.

Розподіл векторів доставки нерівномірний і залежить від типу порушника й цілі. За даними галузевих звітів~\cite{ch01:dbir,enisa} для масових фінансово вмотивованих кампаній домінують фішинг і експлуатація загальновідомих вразливостей, тоді як для цільових операцій проти критичної інфраструктури зростає роль компрометації ланцюга постачання й легітимного віддаленого доступу. Ця відмінність визначає пріоритети захисту для конкретної організації: модель загроз (розділ~\ref{chap:ch04}) має спиратися на релевантний профіль порушника, а не на абстрактний «усереднений» ландшафт. Важливо також, що успішна атака зазвичай поєднує кілька векторів послідовно: початкове проникнення через фішинг, закріплення через легітимні засоби, рух мережею через викрадені облікові дані, --- тому ефективний захист покриває весь ланцюг, а не лише точку входу.

\section{Український контекст}\label{sec:ch01-ukraine}

Україна посідає унікальне місце в історії протидії злоякісному коду як країна, що зазнала найруйнівніших відомих кампаній і водночас накопичила унікальний досвід захисту. \emph{Атаки на енергосистему} (2015, 2016): шкідливий код BlackEnergy і згодом Industroyer~\cite{industroyer} уперше в історії спричинили знеструмлення через кіберзасоби, а Industroyer безпосередньо маніпулював протоколами промислової автоматики (розділ~\ref{chap:ch14}). \emph{NotPetya} (2017): замаскований під вимагач вайпер, доставлений через компрометований механізм оновлення бухгалтерського ПЗ (атака на ланцюг постачання, розділ~\ref{chap:ch15}), завдав глобальних збитків у мільярди доларів, ставши найдорожчою кібератакою в історії~\cite{notpetya,ch01:greenberg}. \emph{Хвиля деструктивних кампаній 2022~року}: з початком повномасштабного вторгнення Україна зазнала безпрецедентної кількості атак вайперами (WhisperGate, HermeticWiper, IsaacWiper, CaddyWiper та інші), спрямованих на знищення даних державних і критичних систем~\cite{wipers2022}, а також оновленого Industroyer2 проти енергетики~\cite{industroyer2}. Аналіз діяльності APT-угруповань проти України систематизовано у вітчизняних дослідженнях~\cite{ch01:apt-ua}.

Ці події дали ключові уроки, що проходять через увесь посібник. По-перше, у воєнному контексті домінують \emph{деструктивні} атаки: мета зловмисника --- не вигода, а знищення даних і порушення функціонування, тому класична логіка «заплати викуп --- віднови дані» непридатна, а єдиним захистом стають ізольовані резервні копії й протестовані плани відновлення (розділ~\ref{chap:ch16}). По-друге, NotPetya продемонстрував критичність ланцюга постачання: компрометація одного довіреного постачальника ПЗ дала зловмиснику доступ до тисяч організацій одночасно (розділ~\ref{chap:ch15}). По-третє, маскування вайпера під вимагач показало, що зовнішня поведінка коду може бути навмисно оманливою, тому атрибуція наміру потребує глибокого аналізу (розділи~\ref{chap:ch06}, \ref{chap:ch07}). По-четверте, швидкість поширення й руйнування (NotPetya паралізував уражені мережі за хвилини) означає, що запобігання й виявлення мають діяти \emph{до} спрацювання корисного навантаження, а реагування --- бути заздалегідь відпрацьованим. По-п'яте, вирішальну роль відіграла національна й міжнародна координація: швидкий обмін індикаторами між CERT-UA~\cite{cert-ua}, Держспецзв'язку~\cite{cip} і світовою спільнотою дослідників дозволив оперативно випустити засоби виявлення й пом'якшення. Статистику й хронологію інцидентів узагальнюють звіти CERT-UA~\cite{ch01:cert-report}, а систематизацію діяльності APT-угруповань --- вітчизняні дослідження~\cite{ch01:apt-ua}. Український досвід перетворив протидію злоякісному коду з технічної на стратегічну дисципліну національної стійкості й зробив вітчизняний науковий доробок у цій сфері затребуваним світовою спільнотою.

\section{Принципи протидії та структура курсу}\label{sec:ch01-structure}

Протидія злоякісному коду ґрунтується на кількох наскрізних принципах. \emph{Глибоко ешелонований захист}: кілька незалежних рубежів, щоб компрометація одного не означала повного ураження. \emph{Припущення компрометації}: виходити з того, що частина атак подолає запобіжні рубежі, і будувати спроможність виявлення й реагування. \emph{Пріоритет виявлення за поведінкою} над сигнатурами, зумовлений масштабом і мінливістю ЗПК~\eqref{eq:ch01-growth}. \emph{Розрив ланцюга ураження} на найраніших досяжних етапах. \emph{Науковість}: оцінювання ефективності засобів за явними метриками з урахуванням базової частоти (розділ~\ref{chap:ch03}).

Структуру курсу побудовано від засад до застосувань. Блок засад (розділи~\ref{chap:ch01}--\ref{chap:ch04}) формує поняття, нормативну базу, теоретичні основи й методологію оцінювання ризиків. Блок аналізу (розділи~\ref{chap:ch05}--\ref{chap:ch09}) розглядає будову ЗПК, статичний і динамічний аналіз, техніки ухилення та машинне навчання. Блок технологій захисту (розділи~\ref{chap:ch10}--\ref{chap:ch11}) --- захист кінцевих точок і мереж. Блок класів загроз і середовищ (розділи~\ref{chap:ch12}--\ref{chap:ch15}) --- вимагачі, APT, спеціалізовані середовища, ланцюг постачання. Блок операційної безпеки й управління (розділи~\ref{chap:ch16}--\ref{chap:ch17}) --- реагування, криміналістика й управління протидією. Методологічною основою слугують класичні праці з будови й аналізу ЗПК~\cite{szor,aycock,sikorski} та інженерії безпеки~\cite{anderson}.

\section*{Висновки до розділу}
\addcontentsline{toc}{section}{Висновки до розділу}

Злоякісний програмний код --- це навмисно створений, несанкціонований код, спрямований на порушення безпеки інформаційної системи; його «злоякісність» визначається наміром і контекстом, а не внутрішніми властивостями, що породжує фундаментальну проблему dual-use і принципову складність автоматичного розрізнення. Таксономія ЗПК є багатовимірною: зразок одночасно класифікується за механізмом поширення (вірус, черв'як, троян), функціональним призначенням (бекдор, вимагач, вайпер, шпигунське ПЗ, руткіт, бот тощо) і способом маскування (пакування, поліморфізм, безфайловість, антианаліз); сучасний код є модульним і багатофункціональним. Еволюція ЗПК відображає зміну мотивації від цікавості до індустрії кіберзлочинності й державних деструктивних операцій, а експоненційне зростання кількості зразків~\eqref{eq:ch01-growth} зробило суто сигнатурний підхід недостатнім. Злоякісний код є інструментом у ширшій операції, життєвий цикл якої описує ланцюг ураження з численними точками виявлення й розриву. Модель порушника варіюється від масових кіберзлочинців до державних APT, для яких відсутні стримувальні чинники~\eqref{eq:ch01-threat}, що вимагає захисту за припущенням компрометації. Україна, зазнавши найруйнівніших кампаній в історії (Industroyer, NotPetya, хвиля вайперів 2022~року), накопичила унікальний досвід, що визначає наскрізні принципи протидії: ешелонованість, припущення компрометації, пріоритет поведінкового виявлення, розрив ланцюга ураження й науковість оцінювання.

\selfcheck
\begin{enumerate}
\item Дайте означення злоякісного програмного коду. Які три ознаки є в ньому ключовими і чому «злоякісність» не визначається внутрішніми властивостями коду?
\item У чому полягає проблема dual-use і які наслідки вона має для автоматичного виявлення ЗПК?
\item Назвіть три осі таксономії злоякісного коду. Чому класифікація є багатовимірною, а не ієрархічною?
\item Порівняйте вірус, черв'як і троян за механізмом поширення й потребою в носії та участі користувача.
\item Опишіть основні класи ЗПК за функціональним призначенням (табл.~\thechapter.1). Чим вайпер відрізняється від вимагача?
\item Які способи маскування злоякісного коду ви знаєте? Як вони пов'язані зі складністю виявлення (розділ~\ref{chap:ch08})?
\item Охарактеризуйте етапи еволюції ЗПК за мотивацією та організацією (рис.~\thechapter.1).
\item За формулою~\eqref{eq:ch01-growth} і співвідношенням $S/\Delta N$ поясніть, чому суто сигнатурний підхід принципово відстає від появи нових зразків.
\item Опишіть етапи життєвого циклу злоякісної кампанії. Чому розрив будь-якого обов'язкового етапу зриває операцію?
\item Порівняйте профілі порушників (табл.~\thechapter.2). Використовуючи~\eqref{eq:ch01-threat}, поясніть особливу небезпеку державних APT.
\item Перелічіть основні вектори доставки ЗПК. Чому зменшення поверхні атаки є необхідним, але недостатнім заходом?
\item Які уроки для протидії злоякісному коду дали атаки Industroyer, NotPetya та хвиля вайперів 2022~року в Україні?
\item Сформулюйте наскрізні принципи протидії злоякісному коду. Чому пріоритет надається виявленню за поведінкою?
\item Дослідницьке: побудуйте фасетну таксономію для вибраного набору реальних зразків ЗПК, визначивши осі класифікації та критерії віднесення, і порівняйте її з категоріями MITRE~ATT\&CK.
\item Дослідницьке: оцініть за даними відкритих звітів динаміку $N(t)$~\eqref{eq:ch01-growth} за останнє десятиліття, визначте коефіцієнт $\alpha$ і проаналізуйте, як він співвідноситься зі спроможністю сигнатурного виявлення.
\end{enumerate}
